\documentclass{article}
\usepackage[UTF8]{ctex}
\usepackage{amsmath,amssymb,amsfonts}
\usepackage{enumitem}
\usepackage{geometry}
\geometry{a4paper, margin=1in}
% === 错题本排版优化（不参与题目标识） ===
\setlength{\parindent}{0pt}
\setlength{\parskip}{0.35em plus 0.10em minus 0.08em}
\setlength{\emergencystretch}{3em}
\setlist{itemsep=0.25em, topsep=0.40em, parsep=0pt}
\allowdisplaybreaks[2]
% === 错题本排版优化结束 ===
\begin{document}
\section*{北京师范大学2025年高等代数真题}
\section*{2．北京师范大学2025年高等代数与空间解析几何试卷}
一．（12 分）设 \(n\) 为正整数且 \(n \geq 2\) ．求

\[
D_{n}=\left|\begin{array}{cccc}
a_{1}+b_{1} & a_{1}+b_{2} & \cdots & a_{1}+b_{n} \\
a_{2}+b_{1} & a_{2}+b_{2} & \cdots & a_{2}+b_{n} \\
\vdots & \vdots & & \vdots \\
a_{n}+b_{1} & a_{n}+b_{2} & \cdots & a_{n}+b_{n}
\end{array}\right|
\]

的值，其中 \(a_{1}, a_{2}, \cdots, a_{n}, b_{1}, b_{2}, \cdots, b_{n}\) 都是实数。

二．（10分）设向量组 \(\alpha_{1}, \alpha_{2}, \alpha_{3}\) 线性无关，向量组 \(\beta_{1}, \beta_{2}, \beta_{3}\) 是

\[
\left\{\begin{array}{l}
\beta_{1}=2 \alpha_{1}+\alpha_{2}+\alpha_{3} \\
\beta_{2}=-\alpha_{1}+3 \alpha_{2}+4 \alpha_{3} \\
\beta_{3}=\alpha_{1}+\alpha_{3}
\end{array}\right.
\]

问 \(\beta_{1}, \beta_{2}, \beta_{3}\) 是线性相关还是无关？并说明理由。

三．（16 分）设

\[
A=\left(\begin{array}{lll}
\lambda & 1 & 0 \\
0 & \lambda & 1 \\
0 & 0 & \lambda
\end{array}\right)
\]

证明：若 \(3 \times 3\) 矩阵 \(B\) 和 \(A\) 乘积可交换，则存在多项式 \(f(x)\) ，使得 \(B=f(A)\) 。

四．（16 分）设

\[
A=\left(\begin{array}{cccc}
1 & -1 & 0 & 2 \\
2 & 0 & 2 & 1 \\
3 & 1 & 4 & 0
\end{array}\right)
\]

求关于变元 \(x_{1}, x_{2}, x_{3}, x_{4}\) 的方程 \(\left(\begin{array}{llll}x_{1} & x_{2} & x_{3} & x_{4}\end{array}\right) A^{\mathrm{T}} A\left(\begin{array}{l}x_{1} \\ x_{2} \\ x_{3} \\ x_{4}\end{array}\right)=0\) 的所有实数解．

五．（16 分）设 \(M_{n}(\mathbb{R})\) 为所有 \(n \times n\) 实矩阵构成的向量空间，设

\[
\begin{gathered}
f: M_{n}(\mathbb{R}) \mapsto M_{n}(\mathbb{R}) \\
X \rightarrow \frac{1}{2}\left(X+X^{\mathrm{T}}\right)
\end{gathered}
\]

证明：\(f\) 为线性变换，并求 \(f\) 的核与秩．

六．（15 分）设 \(A\) 为 \(n\) 阶正定实对称矩阵，\(\alpha, \beta\) 为 \(n\) 维实列向量．证明：

\[
\left(\alpha^{\mathrm{T}} \beta\right)^{2} \leq\left(\alpha^{\mathrm{T}} A \alpha\right)\left(\beta^{\mathrm{T}} A^{-1} \beta\right) .
\]

等号成立当且仅当 \(A \alpha\) 与 \(\beta\) 线性相关．

七．（20 分）填空题，每题5分。\\

\subsection{题目1}
1．\(v=(-1,1,0), w=(1,1,1)\) ，求 \(v \times(v \times(v \times(v \times w)))=\) \(\_\_\_\_\)。\\

\subsection{题目2}
2．直线 \(l_{1}: \frac{x-1}{2}=\frac{y+2}{3}=\frac{2 z-1}{4}, l_{2}: \frac{x+2}{-1}=\frac{y-2}{-2}=\frac{z-\frac{1}{2}}{0}\) ，求 \(l_{1}\) 与 \(l_{2}\) 的距离 \(\_\_\_\_\)．\\

\subsection{题目3}
3．求过直线 \(\left\{\begin{array}{l}x-2 y-z+1=0 \\ 2 x+y=0\end{array}\right.\) 且与平面 \(x-y-z+1=0\) 垂直的平面方程\\
\(\_\_\_\_\)．

八．（12 分）求单叶双曲面 \(\frac{x^{2}}{9}+\frac{y^{2}}{16}-\frac{z^{2}}{4}=1\) 过点 \((3,4,2)\) 的两条直母线方程，并求它们的夹角．

九．（13 分）求过原点且与椭圆柱面 \(\frac{x^{2}}{25}+\frac{y^{2}}{16}=1\) 相交于一个圆周的平面 \(\alpha\) 的方程．

十．（20 分）将 \(3 x^{2}+3 y^{2}+5 z^{2}+2 x y-2 y z-2 x z+2 x+10 y+6 z+7=0\) 化简成标准型，写明使用的坐标变换，并说明是什么二次曲面。
\end{document}
